\section{Introduction}

Vision-language models represent a major advancement in multimodal understanding, achieving state-of-the-art performance on diverse visual reasoning tasks. However, a critical gap remains: can VLMs maintain and update internal state representations across video sequences? This question is essential for applications requiring persistent object tracking, such as embodied AI, activity understanding, and interactive systems.

\subsection{Motivation}

Card-relocation tasks provide an ideal testbed for this question. Unlike natural videos with semantic complexity and causal ambiguity, card-relocation offers:
\begin{itemize}
    \item \textbf{Discrete state space}: positions A/B/C/D (unambiguous ground truth)
    \item \textbf{Observable transitions}: visual card swaps over six rounds
    \item \textbf{Compositional structure}: iterative state updates requiring tracking
    \item \textbf{No extraneous reasoning}: task demands pure state tracking, not world knowledge
\end{itemize}

Trajectory supervision---prompting models to output intermediate states---seems promising, inspired by chain-of-thought prompting's success in language models~\cite{wei2022chain}. Yet a key methodological question remains unanswered: does intermediate supervision improve \emph{genuine} tracking, or primarily teaches models to generate structured outputs matching training distributions?

\subsection{Research Question and Approach}

We ask:
\begin{quote}
\emph{When a VLM is trained to output intermediate visual states (trajectory supervision), does that training repair genuine temporal state tracking, or does it primarily teach output format patterns?}
\end{quote}

To address this, we employ a \textbf{format-matched negative control} (Model B$'$v3): a training dataset identical in structure and statistical properties to ground-truth trajectories, but with labels corrupted to break visual-label correspondence. If trajectory supervision (Model B) does not outperform the corrupted control (B$'$v3), it suggests format learning dominates over semantic tracking.

\subsection{Contributions}

\begin{enumerate}
    \item \textbf{Methodological framework}: A complete pipeline for controlled VLM behavioral studies, including format-matched negative controls and small-sample statistical testing.
    
    \item \textbf{Dataset audit findings}: Discovery of a systematic confound (target\_card $\equiv$ initial\_position) across all 50 videos, demonstrating the necessity of ground-truth validation in VLM evaluation.
    
    \item \textbf{Transparent null reporting}: Honest reporting of a small-sample null result with power analysis and explicit reproducibility materials, rather than p-hacking or post-hoc analysis.
    
    \item \textbf{Complete reproducibility}: 24 journey files documenting experimental iterations, artifact verification tables, and hardcoded hyperparameters enabling reproduction.
\end{enumerate}

\subsection{Scope and Limitations}

This is a \emph{pilot study} with inherent limitations: single 3B model, $n=10$ test videos, synthetic task. We explicitly do not claim trajectory supervision \emph{never} works, only that under these specific conditions (frozen vision tower, 40 training examples, Qwen2.5-VL-3B), it did not show clear advantage over format learning. Larger-scale work with diverse models and tasks is needed.

\subsection{Paper Organization}

We first review related work (\S\ref{sec:related}), then describe the task and experimental design (\S\ref{sec:methods}). Experimental setup and results follow (\S\ref{sec:experiments}, \S\ref{sec:results}). Discussion (\S\ref{sec:analysis}) positions our null result within the broader literature and provides power analysis. Conclusion (\S\ref{sec:conclusion}) outlines future directions.
